\section{Phase 0: plant and fleet validation}
\label{sec:phase0}
Phase~0 verifies the equations, discretization, operating envelope, and physical
family design before identification or control is introduced.

\subsection{Experiment 0A: nominal-vehicle validation}
\paragraph{Objective.}
Validate the nominal continuous and zero-order-hold-discretized bicycle model on
$v_x\in[10,30]\,\mathrm{m\,s^{-1}}$ with initial sample time
$T_s=0.01\,\mathrm{s}$.

\paragraph{Predeclared expectation.}
Matrix entries, poles, time constants, and steady-state gains vary smoothly with
speed. The nominal model is physically plausible and numerically well conditioned
over the complete envelope.

\paragraph{Setup and planned evidence.}
On a fine speed grid, compute continuous and discrete eigenvalues, controllability,
steady-state steering-to-yaw-rate and steering-to-sideslip gains, and an understeer
or equivalent handling measure. Plot selected $A$ and $B$ entries, pole
trajectories, and steady-state gains versus speed. Check equations, units, and sign
conventions independently.

\paragraph{Decision gate.}
No unintended instability, singularity, loss of controllability, discontinuity,
or implausible gain may occur. Any failure blocks all following experiments.

\paragraph{Results, interpretation, and decision.}
Pending.

\subsection{Experiment 0B: structural-family validation}
\paragraph{Objective.}
Verify that the nominal, heavy/high-inertia, and changed-handling-balance centers
create distinct but physically interpretable dynamics.

\paragraph{Predeclared expectation.}
The heavy family differs mainly in response speed and gain. The handling family
changes yaw/sideslip balance rather than merely rescaling the nominal response.

\paragraph{Setup and planned evidence.}
At $v_x\in\{12,20,28\}\,\mathrm{m\,s^{-1}}$, apply a small steering step, a smooth
sine sweep, and a lane-change-like input. Compare yaw rate, sideslip, poles,
steady-state gains, and family-to-family dynamic distance. Report a compact table
of absolute physical parameters and relative changes.

\paragraph{Decision gate.}
All centers must remain controllable and physically plausible. Differences must
be large enough to affect prediction and control, but not so extreme that the
benchmark makes family recovery trivial or unrealistic.

\paragraph{Results, interpretation, and decision.}
Pending.

\subsection{Experiment 0C: within-family client population}
\paragraph{Objective.}
Construct a 30-client fleet with ten clients per family and determine whether the
desired hierarchical structure is actually present.

\paragraph{Predeclared expectation.}
Client envelopes overlap modestly within each family. Persistent differences
between families remain visible over relevant parts of the speed range.

\paragraph{Setup and planned evidence.}
Sample bounded, fixed client parameters with initial ranges
$m,I_z:\pm2.5\%$ and $C_f,C_r:\pm5\%$. Inspect physical parameters, frozen-model
distances, response envelopes, and controllability over speed. The primary figure
shows family centers and within-family bands for selected dynamic characteristics.

\paragraph{Decision gate.}
Within-family variation should be smaller than between-family variation without
perfect separation on every scalar feature. Invalid samples are not silently
discarded; the sampling rule must be revised if they occur systematically.

\paragraph{Results, interpretation, and decision.}
Pending.

\subsection{Phase-0 benchmark freeze}
After Experiments 0A--0C pass, the equations, sample time, speed envelope, family
centers, scatter model, and development seed are frozen. Later changes require an
explicit report entry explaining the failure or new evidence that motivated them.
